
\documentclass{article}
\usepackage{colortbl}
\usepackage{makecell}
\usepackage{multirow}
\usepackage{supertabular}

\begin{document}

\newcounter{utterance}

\centering \large Interaction Transcript for game `cladder', experiment `full\_v1.5\_default', episode 53 with Qwen3.5{-}9B.
\vspace{24pt}

{ \footnotesize  \setcounter{utterance}{1}
\setlength{\tabcolsep}{0pt}
\begin{supertabular}{c@{$\;$}|p{.15\linewidth}@{}p{.15\linewidth}p{.15\linewidth}p{.15\linewidth}p{.15\linewidth}p{.15\linewidth}}
    \# & \multicolumn{2}{c}{Player} && \multicolumn{2}{c}{Game Master} \\
    \hline

    \theutterance \stepcounter{utterance}  
    & & & \multicolumn{4}{p{0.6\linewidth}}{
        \cellcolor[rgb]{0.9,0.9,0.9}{
            \makecell[{{p{\linewidth}}}]{
                \texttt{\tiny{[P1$\langle$GM]}}
                \texttt{You are an expert in causal inference. The following question is not a typical commonsense query, but rather a meticulously designed question created by a professor specializing in causal inference, intended to assess the students' mastery of the course content.} \\
\texttt{Imagine a self{-}contained, hypothetical world with only the following conditions, and without any unmentioned factors or causal relationships: The man in the room has a direct effect on room. The candle has a direct effect on room. Method 1: We look directly at how the man in the room correlates with room in general. Method 2: We look at this correlation case by case according to the candle. To understand how the man in the room affects room, is it more correct to use the Method 1 than Method 2?} \\
\texttt{Start your answer with 'Yes' or 'No', followed by additional reasoning or evidence to support your explanation.} \\
            }
        }
    }
     \\ \\

    \theutterance \stepcounter{utterance}  
    & \multicolumn{4}{p{0.6\linewidth}}{
        \cellcolor[rgb]{0.561,0.737,0.561}{
            \makecell[{{p{\linewidth}}}]{
                \texttt{\tiny{[P1$\rangle$GM]}}
                \texttt{No.} \\
\\ 
\texttt{In the context of causal inference, Method 1 represents an **unadjusted association** (or a naive correlation), while Method 2 represents **stratification** (or conditioning on a variable). To determine the correct approach, we must analyze the causal structure implied by your description:} \\
\\ 
\texttt{1.  **Causal Structure**: You stated that "The man in the room has a direct effect on room" and "The candle has a direct effect on room." This implies a **fork** structure where both the Man (\$M\$) and the Candle (\$C\$) are independent causes of the Room state (\$R\$). There is no mention of the Man affecting the Candle or vice versa, nor is there a confounder affecting both \$M\$ and \$C\$.} \\
\texttt{2.  **The Problem with Method 1**: Looking directly at the correlation between the Man and the Room without considering the Candle (\$P(R|M)\$) yields a biased estimate of the causal effect if the Candle is not independent of the outcome in the observed data distribution. Specifically, if the presence of the Candle varies systematically (e.g., candles are only lit when the man is present, or the candle's effect varies by time of day which also correlates with the man's presence), the marginal correlation will be confounded. Even if \$M\$ and \$C\$ are independent in the population, simply averaging over the distribution of \$C\$ without explicitly accounting for \$C\$'s varying effect on \$R\$ can obscure the specific mechanism or magnitude of \$M\$'s effect if the effect of \$C\$ is heterogeneous.} \\
\texttt{3.  **The Superiority of Method 2**: Method 2 involves looking at the correlation "case by case according to the candle." This is the definition of **stratification**. By conditioning on the Candle (\$P(R|M, C=c)\$), you block any potential backdoor paths (if they existed) and, more importantly in this specific setup, you isolate the effect of the Man within levels of the Candle. This allows you to see if the Man's effect is constant (homogeneous) or if it interacts with the Candle.} \\
\texttt{4.  **Causal Identification**: To correctly identify the causal effect of the Man on the Room, one must generally average the conditional effects (from Method 2) over the distribution of the Candle, rather than relying on the raw marginal correlation (Method 1). Method 1 conflates the direct effect of the Man with the varying influence of the Candle. Therefore, Method 2 is the rigorous, correct approach for estimating the causal effect, whereas Method 1 is prone to bias or misinterpretation.} \\
\\ 
\texttt{Consequently, it is **not** more correct to use Method 1; Method 2 is the necessary step to disentangle the causal influence of the Man from the confounding or modifying influence of the Candle.} \\
            }
        }
    }
    & & \\ \\

    \theutterance \stepcounter{utterance}  
    & & & \multicolumn{2}{p{0.3\linewidth}}{
        \cellcolor[rgb]{0.9,0.9,0.9}{
            \makecell[{{p{\linewidth}}}]{
                \texttt{\tiny{[GM$|$GM]}}
                \texttt{No} \\
            }
        }
    }
    & & \\ \\

    \theutterance \stepcounter{utterance}  
    & & & \multicolumn{2}{p{0.3\linewidth}}{
        \cellcolor[rgb]{0.9,0.9,0.9}{
            \makecell[{{p{\linewidth}}}]{
                \texttt{\tiny{[GM$|$GM]}}
                \texttt{yes} \\
            }
        }
    }
    & & \\ \\

    \theutterance \stepcounter{utterance}  
    & & & \multicolumn{2}{p{0.3\linewidth}}{
        \cellcolor[rgb]{0.9,0.9,0.9}{
            \makecell[{{p{\linewidth}}}]{
                \texttt{\tiny{[GM$|$GM]}}
                \texttt{game\_result = LOSE} \\
            }
        }
    }
    & & \\ \\

\end{supertabular}
}

\end{document}
